% Generated draft: requires mathematical and source review.
\documentclass[11pt]{article}
\usepackage[margin=25mm]{geometry}
\usepackage{fontspec}
\setmainfont{Latin Modern Roman}
\usepackage{amsmath,amssymb,mathtools,unicode-math}
\setmathfont{Latin Modern Math}
\usepackage{enumitem,needspace,xurl,hyperref}
\setlength{\parindent}{0pt}
\setlength{\parskip}{5pt}
\setlength{\emergencystretch}{3em}
\newcommand{\N}{\mathbb{N}}
\newcommand{\Z}{\mathbb{Z}}
\newcommand{\Q}{\mathbb{Q}}
\newcommand{\R}{\mathbb{R}}
\newcommand{\C}{\mathbb{C}}
\newcommand{\F}{\mathbb{F}}
\newcommand{\A}{\mathbb{A}}
\newcommand{\PP}{\mathbb{P}}
\newcommand{\E}{\mathbb{E}}
\newcommand{\eps}{\varepsilon}
\newcommand{\dd}{\,\mathrm{d}}
\begin{document}
\section*{UnsolvedMath ID: 30002110}
\textbf{Source problem code:} catalogue entry
\subsection{Definitions and mathematical statement}
Interpret a bounded planar domain as a nonempty bounded connected open set \(\Omega\subset\R^2\), and write \(A(\Omega)=|\Omega|\) for its Lebesgue area. The record does not specify the boundary conditions or the admissible partition class. The following gives a precise, provisional reading using Dirichlet eigenvalues and the maximum partition energy; these conventions require mathematical review against the source.

For a nonempty bounded open set \(D\subset\R^2\), define its first Dirichlet eigenvalue by
\[
\lambda_1(D)=\inf_{0\ne u\in H_0^1(D)}\frac{\int_D |\nabla u|^2\,\dd x}{\int_D |u|^2\,\dd x},
\]
where \(H_0^1(D)\) is the closure of the smooth compactly supported functions on \(D\) in the Sobolev \(H^1\) norm. For each integer \(k\ge1\), provisionally let
\[
\mathcal P_k(\Omega)=\left\{(D_1,\ldots,D_k):\ D_i\subset\Omega\text{ are nonempty connected open sets, and }D_i\cap D_j=\varnothing\text{ for }i\ne j\right\},
\]
and define
\[
L_k(\Omega)=\inf_{(D_1,\ldots,D_k)\in\mathcal P_k(\Omega)}\max_{1\le i\le k}\lambda_1(D_i).
\]
This provisional admissible class imposes no additional covering or boundary regularity condition. The supplied record does not determine whether such conditions were intended.

Let \(\mathrm{Hexa}_1\) denote the interior of a regular hexagon of area one, and interpret the source notation \(\lambda(\mathrm{Hexa}_1)\) as \(\lambda_1(\mathrm{Hexa}_1)\). With these conventions, the supplied conjecture asks whether, for every bounded planar domain \(\Omega\),
\[
\lim_{k\to\infty}\frac{L_k(\Omega)}{k}=\frac{\lambda_1(\mathrm{Hexa}_1)}{A(\Omega)}.
\]
No boundary regularity hypothesis on \(\Omega\) is supplied. The intended scope of this quantifier and the provisional spectral and partition conventions remain to be checked against the cited report.
\subsection{Short English statement}
The conjecture asks whether the least possible largest first Dirichlet eigenvalue among \(k\) planar partition cells grows like \(k\) times the first eigenvalue of a unit-area regular hexagon, divided by the area of the domain. The Dirichlet and partition conventions used here are provisional because the supplied statement does not define them.
\subsection{Sources}
\begin{itemize}
\item Local UnsolvedMath registry, permanent ID \(30002110\), title ``Hexagonal Asymptotics for Spectral Minimal Partitions,'' source problem code \texttt{catalogue entry}, collection ``Oberwolfach Reports Open Problems.'' The supplied proposed year is 2012; no individual proposer is named.
\item Supplied source citation: \emph{Geometric Aspects of Spectral Theory (2013)}, Oberwolfach Reports, pp. 2013--2076, DOI \texttt{10.4171/owr/2012/33}. \url{https://doi.org/10.4171/owr/2012/33}. The background describes these reports as records of mathematical workshops and open-problem discussions. The registry labels the statement \texttt{corrected\_verified} and reports that the official report confirmed it, with only math delimiters added. This is supplied verification metadata, not an independent verification by this draft.
\item Supplied literature reference: ``Bonnaillie-Noel and Helffer, large k spectral partitions.'' \url{https://arxiv.org/abs/1509.04501}.
\item Supplied literature reference: ``spectral minimal partitions and open problems.'' \url{https://doi.org/10.1007/s00020-021-02635-7}.
\item Supplied literature reference: ``original OWR report.'' \url{https://doi.org/10.4171/owr/2012/33}.
\item Literature assessment dated \texttt{2026-08-21}: the supplied curation labels the problem open and states that lower bounds, numerical evidence, and results for restricted convex or related partition classes do not establish the full spectral asymptotic. This is supplied dated curation, not independently verified. The background explicitly calls for independent verification of time-sensitive status and cited claims.
\end{itemize}
\end{document}
